\section{Experimental development and rejected alternatives}
The project evaluated classical and learned matching pipelines, including DeepMatching, a project-specific gradient-histogram representation, CPM-HOG and modern learned feature matchers, including XFeat (Accelerated Features) \citep{xfeat} and LightGlue \citep{lightglue}. These comparisons showed that network recency or extraction speed alone did not predict DIC coverage. Global affine projection performed well on globally affine tensile tests, but a few informative matches can reconstruct such a field. Those results did not establish direct matching at each requested centre or convergence of the downstream solver.

A second line separated detection from description. DaD (Distilled Reinforcement Learning for Diverse Keypoint Detection) and DeDoDe version 2 \citep{dad,dedode}, together with a lightweight speckle detector, were combined with a common DeDoDe-B description model. DeDoDe is the published model name. This controlled the descriptor while examining point selection. Fast native-resolution detection did not imply a fast complete pipeline: description could dominate pair time. A compact convolutional neural network (CNN) descriptor and learned transformations of a 64-component gradient descriptor were therefore implemented. They established functional training and portable inference, but the compact CNN lost most correspondences in independent rotation, affine and heterogeneous tests. Faster description alone was insufficient.

Affine-compensation experiments subsequently restored many correspondences without additional training. This supported a geometric explanation of failure and motivated explicit orientation hypotheses. Exhaustive orientation banks enabled non-oracle search, but repeated extraction and allocation were costly. The later initialization pipeline retained ten orientations as its coverage reference. An adaptive six-orientation search with two refinements reduced some costs but produced substantial losses on individual difficult pairs and was not adopted as the default. A small average coverage loss is not harmless when concentrated in a region with no remaining seed.

A neural deformation field (NDF) was also evaluated as a residual after affine preconditioning. Its displacement grid was projected to correlation centres through different interpolators. Coordinate conventions required correction before comparisons became interpretable. The corrected experiments still showed that smoothing a wrong displacement grid could not recover the material correspondence. An ablation with and without local normalized cross-correlation (NCC) showed that the residual checkpoint did not improve the best tested chain. That checkpoint had not been trained specifically on post-rectification residuals. The negative result concerns this checkpoint in this role, rather than all learned residual predictors.

The operational conclusion was to retain image-derived geometry and local correlation, confirm seeds with the solver, and propagate accepted local states. This conclusion is based on solved-window coverage. It does not require that the chosen proposal generator be intrinsically superior on every descriptor benchmark. A learned coarse matcher can remain part of hypothesis discovery while a learned residual field is omitted; these are distinct components.

\subsection{CPM-HOG and departures from published CPM}
CPM-HOG is a proposal generator, not a reproduction of the complete published CPM optical-flow pipeline. It uses compact HOG features, sparse-grid matching in both directions, reciprocal checks and local affine projection to requested DIC centres. Published CPM uses dense scale-invariant feature transform (SIFT) Flow features and, in its optical-flow application, an additional dense interpolation stage. Our replacement reduces descriptor storage and permits shared WebGPU kernels. Sparse outputs avoid constructing a dense optical-flow field that the downstream solver does not require.

Parallel propagation changes execution order relative to an in-place sequential scan: an improved state may not influence every later location in the same dispatch. Identical iteration counts need not imply equivalent propagation distances. Local affine projection also changes the output interpretation: a state predicted at a requested centre is supported by measured grid matches, but is not a directly measured correspondence there. The implementation records that distinction.

These modifications trade representational richness and scheduling behaviour for a portable computational path. Historical strong-strain failures prohibit presenting CPM-HOG as the best high-deformation method in the repository. It remains an optional initializer requiring solver validation. More random proposals cannot reliably repair a matching cost that assigns poor similarity to the true deformed support. The propagation policy below therefore remains independent of this particular proposal source.

